\section{L4 级别体验：Deep Research 与 Claude Code}

前面讨论挖矿窗口，本章看当前最明确的矿脉在哪里。广密认为最优秀的两个 Agent 已经有 L4 体验：ChatGPT 的 Deep Research 和 Anthropic 的 Claude Code，分别对应信息搜索/研究和软件开发。今天最大红利仍然是 language/code 红利，尤其是 code，还不是多模态、世界模型或机器人。

\begin{figure}[H]
\centering
\includegraphics[width=0.94\textwidth]{l4-experience-map.png}
\caption{L4 体验地图：Deep Research 和 Claude Code 分别代表搜索和软件开发。自制概念图，依据 00:44:21--00:52:43 对谈内容整理。}
\end{figure}

\begin{knowledgebox}{读图：为什么这两个像 L4}
Deep Research 能完成多步信息搜索、综合和报告；Claude Code 能理解代码库、修改、测试并迭代。这些任务都有明确输入、工具环境和可验证输出，因此更容易形成高阶 Agent 体验。
\end{knowledgebox}

\subsection{为什么 code 红利最大}

代码世界有天然优势：任务可形式化、工具可调用、结果可测试、反馈快、价值高。相比多模态、机器人和世界模型，coding 更早出现 L4 体验并不意外。它是数字世界里最适合 Agent 的高价值环境之一。

\begin{knowledgebox}{L4 体验条件}
\begin{tabularx}{\textwidth}{p{0.22\textwidth}p{0.34\textwidth}X}
\toprule
条件 & Deep Research & Claude Code \\
\midrule
任务清晰 & 研究问题和资料收集 & issue、bug、feature。 \\
工具可用 & 搜索、网页、引用 & 代码库、测试、shell。 \\
结果可评估 & 报告质量和来源 & 测试通过、代码 review。 \\
价值高 & 节省研究时间 & 节省工程时间。 \\
\bottomrule
\end{tabularx}
\end{knowledgebox}

\begin{knowledgebox}{下一批 L4 候选场景}
\begin{tabularx}{\textwidth}{p{0.20\textwidth}p{0.34\textwidth}X}
\toprule
场景 & 为什么可能出现 L4 & 难点 \\
\midrule
法律/合规 & 文档密集、流程明确、价值高 & 责任重，错误成本高。 \\
金融研究 & 信息检索、模型、报告和交易线索 & 数据授权和实时性要求高。 \\
销售运营 & CRM、邮件、会议、报价都可工具化 & 私有数据和组织流程复杂。 \\
数据分析 & SQL、报表、BI、业务解释闭环 & 需要理解指标口径和业务语义。 \\
设计/多模态 & 视觉生成和编辑迭代变快 & 主观评价和版权问题更复杂。 \\
\bottomrule
\end{tabularx}
\end{knowledgebox}

\subsection{本章小结}

L4 体验不是抽象等级，而是用户真的愿意把复杂任务交给系统。Deep Research 和 Claude Code 是目前最清晰的两个样本。

